%! Project: Design and Conduct a Survey — Unit 1, Lesson 1.8 — Guided Notes (Answer Key)
\documentclass[10pt]{article}
\usepackage{saar-article}
\usepackage{saar-key}   % pulls in -boxes; do NOT also load -boxes
\newcommand{\TallMath}[1]{$\displaystyle #1\rule[-1.4em]{0pt}{3.2em}$}

\begin{document}

\pageheader{Unit 1, Lesson 1.8}{Guided Notes --- Answer Key}
% NAMESTRIP: no \namedateperiod here — mirrors the blank. Teacher-only prose goes
% in the lesson plan as \begin{teachernote}[Guided Notes], never in a key.

\vspace{0.15in}

\begin{objectivebox}
\small
By the end of this lesson, I will be able to\dots
\begin{enumerate}[leftmargin=1.5em, itemsep=2pt, topsep=2pt]
  \item plan a survey through all four stages of the \textbf{statistical
        cycle};
  \item build a \textbf{survey instrument} whose items are neutral and whose
        choices can be counted;
  \item \textbf{collect} real data from a real sample and organize it into a
        frequency table;
  \item choose and complete a \textbf{visual representation} that fits
        categorical data;
  \item report a conclusion \emph{in context} and judge its \textbf{validity}
        --- what my design does and does not let me claim.
\end{enumerate}
\end{objectivebox}

\boxguard
\begin{vocabbox}
\small
Fill in each term as we build it together.
\vocabans{Survey instrument}{The actual list of questions, in the order they
  will be asked, that every person in the sample answers. It is the tool that
  turns a statistical question into data.}
\vocabans{Closed-response item}{A question offering a fixed set of choices, so
  every answer falls into one category and the answers can be counted. An
  open-response item is answered in the person's own words instead.}
\vocabans{Relative frequency}{The count in one category divided by the total
  number of responses, written as a decimal or a percent. The relative
  frequencies of all the categories add to $1$, or $100\%$.}
\vocabans{Bar graph}{A display for categorical data in which the height of each
  bar is the frequency or the relative frequency of one category. The bars are
  drawn separated, and the categories may be placed in any order.}
\vocabans{Validity of a conclusion}{Whether the study's design actually supports
  the claim being made --- who was sampled, how they were chosen, and therefore
  whom the result is allowed to describe.}
\end{vocabbox}

\boxguard[12]
\begin{hookbox}
\small
The Student Council has \$$3{,}000$ and two members are arguing. One says
\emph{``everybody wants longer lunch.''} The other says \emph{``nobody I know
does.''} Both of them are describing the people they happen to sit with.

\vspace{2pt}
The council's stratified sample of $60$ students says something neither of them
said: club funding $35\%$, field trips $30\%$, gym equipment $25\%$, longer
lunch $10\%$.

\vspace{2pt}
\textbf{What does the council have that the two arguing members do not? Name it
in one word, and then say what it cost them to get it.}
\ansline{Data --- and it cost them a plan: naming the population, choosing $60$ students by chance, and asking all of them the same question.}
\end{hookbox}

\boxguard[24]
\begin{notesbox}{1. Your Project, on One Page}
\small
Today you run a study of your own. Nobody is assigned to anything --- you ask
and you record --- so what you are running is an \textbf{observational study},
and the plan is the same four-stage cycle you have used since Lesson 1.1.

\begin{center}
\renewcommand{\arraystretch}{1.4}
\begin{tabularx}{0.98\linewidth}{@{} p{2.7cm} p{5.3cm} X @{}}
\toprule
\textbf{Stage} & \textbf{What the stage asks} & \textbf{What you hand in} \\
\midrule
Formulate & What do we want to know, and about whom?
  & \ans{One statistical question and the population it is about.} \\
Collect & Who do we ask, how does chance choose them, and how do the answers
  arrive? & \ans{A sampling plan and the finished instrument.} \\
Represent & How do we organize what came back?
  & \ans{A frequency table and a bar graph.} \\
Analyze \& report & What do the numbers say \emph{in context} --- and what do
  they \emph{not} say? & \ans{A written report with its limits named.} \\
\bottomrule
\end{tabularx}
\end{center}

\vspace{2pt}
\textit{Every stage hands something to the next one. A bad question cannot be
repaired by good arithmetic later.}
\end{notesbox}

\boxguard[26]
\begin{notesbox}{2. Stage 1 and Stage 2 --- The Question and the Sample}
\small
A \textbf{statistical question} passes three tests: the answers \emph{vary}, a
\emph{group} is named, and a \emph{measurement} is named. Check each candidate.

\begin{center}
\renewcommand{\arraystretch}{1.4}
\begin{tabularx}{0.98\linewidth}{@{} X c p{4.3cm} @{}}
\toprule
\textbf{Candidate question} & \textbf{Passes?} & \textbf{If not, repair it} \\
\midrule
How should the council spend the grant? & \ans{No} & \ans{No group and no
  measurement named --- it asks for an opinion, not data.} \\
Which of four options would the most Riverbend students choose for the grant?
  & \ans{Yes} & \ans{---} \\
\bottomrule
\end{tabularx}
\end{center}

\vspace{3pt}
Now stage 2. The council wants $60$ answers from $900$ students and wants every
grade represented, so it uses a \textbf{stratified} sample at the rate
$60 \div 900 = \frac{1}{15}$.

\vspace{2pt}
How many $9$th graders (there are $270$) does that rate call for?

\begin{work}
  270 \times \dfrac{1}{15} &= 18 \text{ students}
\end{work}

How many $12$th graders (there are $180$)?

\begin{work}
  180 \times \dfrac{1}{15} &= 12 \text{ students}
\end{work}

\vspace{2pt}
The $240$ tenth graders give $16$ and the $210$ eleventh graders give $14$, so
the four allocations total \ans{$60$}. Finally, the answers have to
\emph{arrive} somehow: of the five methods from Lesson 1.7, the one that fits
``a fixed question answered by many people'' is a \ans{survey}.

\vspace{2pt}
\textbf{Careful --- this is about you.} When your class surveys only the
students in this room, chance did not choose them; the schedule did. That makes
your sample a \ans{convenience sample}, and you will owe your reader a weaker
sentence because of it.
\end{notesbox}

\boxguard[26]
\begin{notesbox}{3. Building the Instrument}
\small
The \textbf{instrument} is the actual list of questions you will read out, word
for word, to every person. Four rules:

\begin{enumerate}[leftmargin=1.6em, itemsep=1pt, topsep=3pt]
  \item \textbf{One idea per item.} Two ideas in one item cannot be answered by
        one person.
  \item \textbf{Neutral wording.} The item must not tell the person which
        answer you want.
  \item \textbf{Choices that cover everyone and do not overlap.} Then the
        percents add to $100\%$.
  \item \textbf{Short.} A person answers what they finish reading.
\end{enumerate}

\vspace{3pt}
Each item below breaks one rule. Name the problem --- you met all three of these
as \textbf{response bias} in Lesson 1.4 --- and repair the item.

\begin{center}
\renewcommand{\arraystretch}{1.45}
\begin{tabularx}{0.98\linewidth}{@{} X p{2.8cm} p{4cm} @{}}
\toprule
\textbf{The item as written} & \textbf{Which rule?} & \textbf{Your repair} \\
\midrule
``Don't you agree that our underfunded clubs deserve the grant?''
  & \ans{Rule 2 --- leading} & \ans{``Which one option should the grant
  fund?''} \\
``Should the council fund clubs and field trips?''
  & \ans{Rule 1 --- two ideas} & \ans{Ask about clubs and field trips as
  separate choices.} \\
``How old are you: $14$--$16$, $16$--$18$, or $18+$?''
  & \ans{Rule 3 --- overlap} & \ans{``$14$--$15$, $16$--$17$, $18$ or
  older.''} \\
\bottomrule
\end{tabularx}
\end{center}

\vspace{2pt}
\textit{A biased instrument produces data that are wrong in a direction. No
amount of arithmetic afterwards can pull them back.}
\end{notesbox}

\boxguard[30]
\begin{notesbox}{4. Represent, Then Analyze}
\small
The council's $60$ answers came back. Complete the \textbf{relative frequency}
column --- each count divided by $60$, written as a percent.

\begin{center}
\renewcommand{\arraystretch}{1.35}
\begin{tabularx}{0.86\linewidth}{@{} X c c @{}}
\toprule
\textbf{Choice} & \textbf{Frequency} & \textbf{Relative frequency} \\
\midrule
Club funding   & $21$ & \ans{$35\%$} \\
Field trips    & $18$ & \ans{$30\%$} \\
Gym equipment  & $15$ & \ans{$25\%$} \\
Longer lunch   & $6$  & \ans{$10\%$} \\
\midrule
\textbf{Total} & $60$ & \ans{$100\%$} \\
\bottomrule
\end{tabularx}
\end{center}

\vspace{3pt}
The variable ``choice'' is \ans{categorical}, so the display that fits is a
\textbf{bar graph}: one bar per category, bars separated, height $=$ relative
frequency. The axes are drawn and scaled for you --- draw the four bars.

\begin{center}
\begin{tikzpicture}[x=1cm, y=0.088cm]
  \foreach \y in {5,15,25,35}
    \draw[linegray, very thin] (0,\y) -- (10.4,\y);
  \foreach \y in {10,20,30,40}
    \draw[linegray] (0,\y) -- (10.4,\y);
  \foreach \y in {0,10,20,30,40}
    \node[left, font=\scriptsize] at (-0.05,\y) {\y\%};
  \foreach \x in {2.6,5.2,7.8}
    \draw[linegray, dashed] (\x,0) -- (\x,42);
  \foreach \i/\lab in {1/Club funding, 2/Field trips, 3/Gym equipment,
                       4/Longer lunch}
    \node[below, font=\scriptsize, align=center, text width=2.4cm]
      at ({2.6*\i-1.3},-0.5) {\lab};
  \draw[cerulean, thick] (0,42) -- (0,0) -- (10.4,0);
  % Key only: the four correct bars, drawn to the heights answered above.
  \foreach \i/\h in {1/35, 2/30, 3/25, 4/10}
    \fill[keyred, opacity=0.55] ({2.6*\i-2.0},0) rectangle ({2.6*\i-0.6},\h);
\end{tikzpicture}
\end{center}

\vspace{4pt}
Finally, \textbf{analyze}: about how many of all $900$ Riverbend students would
choose gym equipment?

\begin{work}
  0.25 \times 900 &= 225 \text{ students}
\end{work}
\end{notesbox}

\boxguard
\begin{practicebox}
\small
\textbf{Stage 4 --- say it, then check that you are allowed to say it.} A report
is four sentences. Write them for the council's survey.

\begin{enumerate}[leftmargin=1.6em, itemsep=3pt, topsep=3pt]
  \item \textbf{What we asked, and of whom.}
        \ansline{We asked which of four options Riverbend's $900$ students would most want the \$$3{,}000$ grant to fund.}
  \item \textbf{How the sample was chosen, and how many answered.}
        \ansline{A stratified sample of $60$ students, $18$/$16$/$14$/$12$ by grade; all $60$ answered.}
  \item \textbf{What we found} --- the two largest categories, as percents.
        \ansline{Club funding led at $35\%$, followed by field trips at $30\%$.}
  \item \textbf{What it means for the \$$3{,}000$}, and about whom.
        \ansline{Club funding is the single choice most Riverbend students support --- about $315$ of the $900$.}
\end{enumerate}

\vspace{2pt}
Now the \textbf{validity} check. The council may write sentence 4 about all
$900$ students. Name the one feature of its design that earns it that right,
and say what would have happened to that right if it had surveyed only the
$60$ students in the cafeteria at noon.
\ansline{Chance chose the sample, and every grade was represented in proportion. A noon cafeteria group is a convenience sample, so the result would describe only those $60$ students --- and a longer-lunch question asked at lunch invites bias on top of that.}
\end{practicebox}

\end{document}
